\chapter{Open assumptions}
\label{app:assumptions}

This appendix collects the quantities that are assumed rather than measured
or read from a datasheet. They are gathered in one place because these are
the points at which the design would change if the assumption turned out to
be wrong.

They are ordered by how much they would change the result. Their numbers are
identifiers rather than positions: the chapters and the bench plan refer to
them, so they were kept when an item moved between sections.

\section{Would change the design}

\paragraph{A.1 Pollution degree 1, obtained by conformal coating.}
Every insulation distance in this design, functional and reinforced
alike, is qualified in pollution degree 1, which is claimed on the basis
that the high-voltage area is conformally coated on both faces. The
coating is an existence condition: in pollution degree 2
the functional requirement at the \SI{900}{\volt} working voltage is
\SI{9.00}{\milli\metre} on ordinary FR-4 where the \device{} package offers
\SI{3.24}{\milli\metre} between its drain and source pads, and no layout can
recover a factor of \num{2.8}. At \SI{900}{\volt} the coating also carries
the reinforced barrier under the isolator bodies (A.24). If the coating is refused, the board is not buildable with
this device, and the answer is a different package, not a different layout.

The assumption is a \emph{process} condition rather than a dimension. A design rule check verifies
dimensions and passes the board either way. This is why it is written at the
head of the rules file, on the fabrication drawing, and here. See
\cref{sec:insulation}.

It is no longer only an assumption of process. The coating is specified with
a UV tracer and every board is inspected under UV light over the whole
high-voltage area (\cref{sec:isoslot}), so each board carries the evidence
that the condition holds on it. What remains assumed is the coating's
long-term integrity, not its presence.

\paragraph{A.2 Rated impulse withstand voltage, and the non-applicability of
overvoltage category.}
Clearance follows from an assumed rated impulse voltage, and because creepage
may not be smaller than clearance, that assumption can also set the creepage.
\SI{2.5}{\kilo\volt} is declared, derived from the transient the converter
generates internally (\SI{1148}{\volt} at the die in the worst simulated
case) rather
than from an overvoltage category, on the grounds that a battery-fed floating
system is not supplied from a low-voltage distribution network and cannot
receive a mains-borne impulse (\cref{sec:ovc}).

This is the critical entry of this group. The design passes at \SI{4}{\kilo\volt} with
\SI{8}{\percent} of margin and fails at \SI{6}{\kilo\volt}, which is what
overvoltage category II on a \SI{1000}{\volt} band would give. The argument
rests on three legs, and it falls if any one of them falls: the standard
provides for deriving the impulse from the actual transient where the circuit
is not mains-supplied; that transient is simulated rather than assumed; and
the declared level is \num{2.2} times the worst simulated peak.

\paragraph{A.3 Table values of IEC 60664-1, and their source.}
The creepage and clearance requirements used throughout
\cref{sec:insulation} come from a third-party calculator implementing a
stored subset of the standard~\cite{standardclarity}, not from the text of
IEC 60664-1~\cite{iec60664}, which was not available.

The calculator checked nine values that had been quoted from memory (five
creepage rows, two clearance rows, the reinforced doubling rule and the
pollution-degree floor), and all nine agreed to the digit. A published,
reproducible secondary source that confirms every value is a different kind
of evidence from a recollection. It is still not the standard: it covers a
stored subset and describes itself as a reference estimate, so the text is to
be read before release. A.2 weighs more than this entry.

\paragraph{A.4 Galvanic isolation in the charging path.}
The argument that overvoltage category does not apply (A.2) depends on the
high-voltage bus never being conductively tied to the mains. During charging
the contactors are closed and the bus \emph{is} connected to the inlet, so
what breaks the path is the galvanic isolation of the charging
equipment (the on-board charger's transformer in AC charging, the station's
in DC charging). Both sit outside this perimeter, and the assumption is
recorded as requirement R10 (\cref{tab:sysreq}).

It is not normatively guaranteed. IEC 61851-23~\cite{iec6185123}
distinguishes isolated from non-isolated DC charging stations, imposing
additional protective measures on the latter, so a vehicle that accepts any
station cannot assume isolation. If the assumption fails, the insulation
coordination moves into a different normative domain rather than
losing margin, and the mitigation belongs at the charge inlet, not on this
board. This is the second most load-bearing assumption in the
design, after A.1.

\paragraph{A.5 Short-circuit withstand time, carried as a threshold.}
The \device{} datasheet does not declare $t_\mathrm{SC}$, the only
selection criterion the chosen device leaves undemonstrated. It is carried
as a threshold (\cref{sec:tsc}): counted as the manufacturer measures it,
the protection chain closes in \SI{0.69}{\micro\second} typical and
\SI{0.83}{\micro\second} worst case, so the design survives a short
if $t_\mathrm{SC} \geq \SI{0.83}{\micro\second}$ and has a factor of two at
\SI{1.66}{\micro\second}. A class value from the manufacturer, on a
\SI{16}{\milli\ohm} die of the same generation at \SI{+15}{\volt}, gives
\SI{2.4}{\micro\second} at \SI{950}{\volt} once scaled at constant
energy~\cite{wolfspeed_scan}: a factor of 2.9, but on another part. The
plan B is sized (\cref{sec:planb}): with a \SI{4.7}{\pico\farad} blanking
capacitor and a \SI{4.7}{\volt} zener the protection covers
\SI{0.71}{\micro\second}; below that, desaturation detection is the wrong
principle for this device. The value is to be obtained from the
manufacturer in writing or from test B1 of \cref{tab:benchplan}, before the
protection is trusted.

\paragraph{A.6 Switching energy: model against datasheet.}
At the datasheet test point the SPICE model gives $\Eon = \SI{748}{\micro\joule}$
against \SI{2552}{\micro\joule} measured, a factor of 3.4, and
$\Eoff = \SI{509}{}$ against \SI{788}{\micro\joule}, a factor of 1.55
(\cref{sec:kelvin}); the device of the first iterations showed the same
factor on turn-on. The loss budget uses the datasheet, on the grounds that it
is a measurement, but this is the first thing to verify on the
bench: if the model is right, the thermal budget improves considerably. The
overshoot results do not depend on it: they come from the same model's
waveforms, not from its energies.

\paragraph{A.7 Insulating pad thermal resistance.}
$\Rthcs = \SI{0.50}{\celsius\per\watt}$ is a class value, not a product
figure, and it also depends on the mounting torque. Across the plausible
range it moves the nominal junction temperature by \SI{23}{\celsius}
(\cref{tab:rthcs}) and never threatens the rated point. At locked rotor it
does: C1 holds only up to \SI{0.62}{\celsius\per\watt}, about
\SI{7}{\celsius} per tenth (\cref{sec:lockedrotor}). The pad value is
therefore a mounting specification, torque included, and not a free choice;
it is the largest single uncertainty in the numbers quoted and the first
thing to measure on a mounted sample. On the real rail (\cref{sec:cooling})
the limit tightens to \SI{0.50}{\celsius\per\watt}, and the design adopts
aluminium nitride with a phase-change interface material at
\SI{0.18}{\celsius\per\watt}; the interfaces decide the value more than the
ceramic does (\cref{sec:coolalt}).

\paragraph{A.8 Sampling rate and sampling instant.}
The analogue interface of \cref{sec:analogif} is designed for
$f_s = \fsw = \SI{20}{\kilo\hertz}$, one sample per switching period, taken
synchronously with the carrier peak or valley. This is standard practice in
machine control and it is what makes a current loop of a few kilohertz
sensible, but the controller has not been selected and the rate has not been
agreed with anyone.

Two things follow from it. The anti-aliasing corner is placed at
\SI{106}{\kilo\hertz} rather than near Nyquist \emph{because} synchronous
sampling already nulls the ripple; if the sampling turned out to be
asynchronous, the ripple would alias, the filter would have to move down by
a decade, and the \SI{16}{\micro\second} of delay that follows would not fit
the current loop, and the board would then need its own ADC, which reverses the
boundary decision of \cref{sec:sysreq}. And the reading error moves from zero
to $\pm\SI{4.10}{\percent}$ of rating (\cref{tab:sampling}).

It is therefore written as requirement R6 rather than left as an
assumption, because the board depends on something that happens outside
it. It is listed here because a requirement that the
system cannot meet is an assumption again.

\paragraph{A.22 Coolant temperature.}
The \SI{65}{\celsius} coolant is this design's assumption, not part of the
specification (\cref{tab:spec}). At \SI{70}{\celsius} the locked-rotor
junction temperature rises by about \SI{6}{\celsius}: on the ideal plate of
\cref{sec:lockedrotor} it stays within C1 with a few degrees to spare, on the
cooling rail adopted in \cref{sec:coolalt} it reaches \SI{117.8}{\celsius} only
with every other tolerance stacked on top (\cref{tab:coolalt}). It is a requirement on the vehicle's cooling loop and should be
written as one.

\paragraph{A.23 Maximum continuous battery voltage.}
Every working-voltage quantity (creepage, the static derating of the
devices, the voltage class of the film capacitor, the divider) is sized at
\SI{900}{\volt}. That is requirement R12 on the battery, including charging
and regenerative braking; a pack that sits above it for long enough to count
as ``continuous'' is outside this design.

\paragraph{A.25 The isolated DC/DC converters (open decision).}
The eleven converters that supply the drivers and the sensing cross the
HV-to-control barrier, and their datasheet recognises them as reinforced
insulation only up to \SI{150}{\volt} rms, basic up to \SI{300}{\volt} rms,
functional up to \SI{2.4}{\kilo\volt}~\cite{murata_mgj2}. At
\SI{900}{\volt} they are functional insulation. The reinforced
barrier declared in this report is therefore not met at these components.
Two ways out, both changing the design: supplies rated for reinforced
insulation at a working voltage of at least \SI{900}{\volt} DC, or a
safety barrier placed elsewhere in the system with the board's HV-to-control
insulation declared functional. It is the one open point of this appendix
that is a safety question rather than a margin question.

The schematic and the board were frozen with these converters, so
the frozen board implements the second option as far as the board can: the
rules file treats the converter pins as functional insulation, at
\SI{5.5}{\milli\metre}, and declares it by name (\cref{tab:rules}); every
other crossing of the barrier (isolators, board copper, the slots) is
reinforced. Choosing the first option means a different module and a
different footprint, that is a board revision, not an edit.

\section{Would change a component choice}

\paragraph{A.9 Negative gate bias (resolved).}
\SI{-5}{\volt} is used. The datasheet's \SI{-4}{\volt} is a typical
operating value, not a limit; the absolute maximum is
\SI{-8}{\volt}~\cite{wolfspeed_e3m0016}, leaving \SI{3}{\volt} of margin. The
standard isolated module delivers \SI{-5}{\volt} without an extra
regulator, and the extra volt goes straight into the margin against spurious
turn-on (\cref{tab:miller}). See \cref{sec:negbias}.

\paragraph{A.10 Isolation grade of the gate-drive supply.}
The \dcdc{} withstands \SI{2400}{\volt} continuously but its \emph{certified}
reinforced rating is only \SI{150}{\Vrms}~\cite{murata_mgj2}. Functional
isolation has been assumed, on the basis that the HV system floats with
respect to chassis and the low-voltage side is not touchable. If certified
reinforced isolation is required across that barrier, the module must change.

\paragraph{A.11 Vehicle \SI{12}{\volt} rail range.}
Assumed \SI{9}{}--\SI{16}{\volt} continuous with larger transients, typical
of ISO 16750-2~\cite{iso16750}, without access to the standard itself. The
pre-regulator must be sized on the real range.

\paragraph{A.12 Terminal dimensions and voltage class of the ceramic capacitors.}
The decoupling capacitor carries up to \SI{900}{\volt} between its own two
terminals, so its creepage is the body length less the two terminations. On
the typical dimensions used here a 1210 leaves \SI{2.00}{\milli\metre}, below
the \SI{2.80}{\milli\metre} required, while a 1812 leaves
\SI{3.20}{\milli\metre}. The voltage class is \SI{1250}{\volt} or more, in
C0G: a \SI{1000}{\volt} part would work at \SI{90}{\percent} of its rating
at \SI{900}{\volt}, and an X7R one would lose most of its capacitance to the
bias. Resolved at capture (D8, \cref{tab:captdecisions}): a
\SI{220}{\nano\farad} part of that class does not exist; the largest C0G
ceramic at \SI{1200}{\volt} is \SI{47}{\nano\farad} in a 3640, whose
\SI{6.25}{\milli\metre} between pads passes with margin. Four are
fitted per leg, \SI{188}{\nano\farad}, at \SI{75}{\percent} of the rating at
\SI{900}{\volt}~\cite{kemet_kclink}.

\paragraph{A.13 Material group of the ceramic pad.}
The pad is sized in pollution degree 2 with the ceramic assigned material
group I, which requires a datasheet stating CTI $\geq \SI{600}{\volt}$.
Technical ceramics (aluminium nitride, adopted in \cref{sec:coolalt}, and
alumina) do not track, so the assignment is expected to hold, but it has not
been confirmed on a specific part. Material group II would take a pad
standing on its own from \SI{33}{}$\times$\SI{37}{} to
\SI{41}{}$\times$\SI{45}{\milli\metre} (\cref{tab:pad}); on the cooling rail,
where the creepage below the tab comes from raising the rail, it would mean
raising it further and checking the section of \cref{fig:mechsection} again.

The pollution degree at the pad is not in this appendix. Degree 1 at the pad
would rest on a sealed housing, a system-level property this design does not
own. Sizing the pad for
degree 2 removed the assumption at a cost of \SI{57}{\percent} more pad
area, the only point in this design where an assumption was eliminated
rather than declared.

\section{Would change the layout or the verification}

\paragraph{A.14 Loop inductance of the real layout.}
The \SI{20}{\nano\henry} limit is not yet a verified result, and it is now
load-bearing: $R_{g,\mathrm{off}} = \SI{12.5}{\ohm}$ keeps the die below
\SI{1080}{\volt} \emph{at} \SI{20}{\nano\henry}, with six and a half volts to spare.
In inductance that margin is \SI{2.6}{\nano\henry}: C2 fails at
\SI{22.6}{\nano\henry} at turn-off and at \SI{23.6}{\nano\henry} on
recovery (\cref{tab:lthreshold}), each known to about
$\pm\SI{1}{\nano\henry}$. Extracted on the frozen board
(\cref{sec:loopextract}), it is \SI{7.8}{\nano\henry} of copper in revision 1.1,
about \SI{8.3}{\nano\henry} with the capacitor ESL, so the verification
holds with \SI{14}{\nano\henry} to the threshold, and the double pulse with
the extracted values puts the die at \SI{1040}{\volt}. What remains assumed is
the extraction itself (a quasi-static model of the copper layers, with the
packages taken from the vendor's model) until test B2 measures it,
against the failure threshold rather than the limit, and
remembering that the figure to compare is the PCB contribution and not the
whole loop. Between \SI{20}{} and \SI{22.6}{\nano\henry} the design holds
as it is; above \SI{23}{\nano\henry} both gate resistors, the losses and the
thermal result of \cref{sec:lockedrotor} have to be recomputed.

\paragraph{A.15 IPC-2152 chart values.}
The conductor widths use the IPC-2221B closed form, which is conservative in
this configuration. The IPC-2152 charts were not available and their values
have not been invented. The layout was frozen without them; the one place
where the pours were narrower than the closed form asks, a neck below each
shunt pad, has a second pour on L6 since revision 1.1
(\cref{sec:conductors}), so the charts are expected to confirm rather than
change the result.

\paragraph{A.16 Copper temperature rise.}
A \SI{30}{\celsius} rise was assumed for conductor sizing
(\cref{tab:widths}). On a board already at \SI{67}{}--\SI{85}{\celsius} this
puts the copper near \SI{115}{\celsius}, which should be revisited against the
mechanical thermal design. The margin here is thinner than it was: the DC bus
current is \SI{32.0}{\ampere} rather than the \SI{22.5}{\ampere} first
assumed, so the bus conductor has to be \SI{60}{\percent} wider.

\paragraph{A.17 Dielectric strength of the prepreg, and the L1--L2 field.}
The \SI{0.10}{\milli\metre} prepreg chosen to minimise the plane-pair
inductance carries the full DC-link voltage, a field of
\SI{9}{\kilo\volt\per\milli\metre} at \SI{900}{\volt} against a laminate strength assumed to be
\SI{30}{}--\SI{40}{\kilo\volt\per\milli\metre}. This is functional
insulation, so the standard imposes no minimum thickness, but it is the
thinnest margin in the design and it cannot be inspected once the board is
pressed. A layer-to-layer withstand test at \SI{2}{\kilo\volt} DC for one
minute is specified on the fabrication drawing, which converts the assumption
into a measurement. Doubling the prepreg would halve the field and
cost \SI{0.19}{\nano\henry} (\cref{tab:planepair}), a trade to revisit at the
first prototype.

\paragraph{A.18 Y-capacitance permitted by the vehicle.}
The board carries \num{2}$\times$\SI{4.7}{\nano\farad}, not the
\num{2}$\times$\SI{22}{\nano\farad} first assumed: \SI{4.7}{\nano\farad} is
the largest Y1 disc rated \SI{1500}{\volt} DC (D12,
\cref{tab:captdecisions}). The function (closing the common-mode current to
chassis at high frequency) stays; its effectiveness is a matter for the
emission test, and the value must still be reconciled with the vehicle's
insulation-monitoring budget, which is a system-level constraint and not a
board-level one.

\paragraph{A.19 CISPR 25 edition and limit class.}
The estimate of \cref{sec:emcest} shows that an input filter is needed
whatever the class; the class decides its margin. It was taken against the
class~3 voltage-method limits of the low-voltage network as a reference,
since the limits that apply to the high-voltage lines and the class the
vehicle maker asks for are not fixed.

\paragraph{A.24 The reinforced barrier under the isolator bodies (decided: slot).}
A through slot is made under each of the seven drivers and four isolated
amplifiers (\cref{tab:slotdims}), and the driver footprint is changed to the
manufacturer's high-voltage land pattern, because the standard one left
\SI{7.25}{\milli\metre} and not the \SI{8}{\milli\metre} assumed until then.
The board path around each slot is at least \SI{9.8}{\milli\metre} against the
\SI{9.5}{\milli\metre} rule. What stays assumed is the package
surface: \SI{8.0}{} and \SI{8.5}{\milli\metre} against \SI{9.0}{\milli\metre}
in pollution degree 2, covered by the components' own reinforced rating and,
in the design basis, by the coating. The slot moves the barrier's weak point
from the uninspectable gap under the body to package faces that are coated
and inspected (\cref{sec:isoslot}).

\section{Accepted limitations}

These are known limits of the design as scoped, collected here so that
nobody has to discover them by reading every chapter.

\paragraph{A.L1 The active discharge depends on the auxiliary supply.}
The seventh gate-drive channel that switches the discharge resistor is fed
from the vehicle's \SI{12}{\volt} rail. If a crash interrupts that rail the
active discharge does not run and only the passive bleeder remains, which
takes \SI{121}{\second} to reach \SI{60}{\volt} from \SI{900}{\volt} against a
\SI{5}{\second} requirement~\cite{unr100}. This is the standard
architecture, in which the bleeder is a backstop and not a substitute, but
the gap is
recorded because it is a requirement unmet in that condition. A production
design would take the discharge control power from the HV bus itself.
See \cref{sec:dclinkmgmt}.

\paragraph{A.L2 Contactors, precharge and interlock monitoring are out of
scope.} The board receives an already-charged bus and carries the interlock
loop through without monitoring it. Its contribution is the DC-link voltage
measurement on which the vehicle controller closes the main contactor at
voltage rather than on a timer. The precharge sizing is given in
\cref{sec:scope} as a system-level result, not as a board function.

\section{Would change only the transient}

\paragraph{A.20 Cold-plate heat capacity and thermal resistance.}
$C_\mathrm{hs} = \SI{1000}{\joule\per\kelvin}$ is estimated from an assumed
plate geometry; it affects no steady-state temperature (verified by
re-running the complete sweep), but it sets the time constant, and therefore
how long a bench measurement must run before the reading is meaningful.
$\Rthha = \SI{0.01}{\celsius\per\watt}$ is no longer an assumption of the
design: it was the optimistic end of what cold plates achieve, and
\cref{sec:cooling} replaces it with the rail as drawn, computed in three
dimensions (equivalent to \SI{0.069}{\celsius\per\watt}). What remains
assumed there is the channel correlation, whose $\pm\SI{20}{\percent}$ moves
the hottest junction by \SI{0.8}{\celsius}, and a uniform flux under the
tab.

\section{Would change the safe state}

\paragraph{A.21 Load type.}
The entire report assumes a passive $RL$ bench load. With a permanent-magnet
machine the safe state would have to be re-examined: the back-EMF can feed the
DC-link with all devices off, and the usual answer is an active short circuit,
which is the opposite of turning everything off.

\section*{Bench verification plan}
\addcontentsline{toc}{section}{Bench verification plan}

The thresholds of this appendix are useful only if each is attached to a
measurement with a pass mark written before the measurement is made.

\begin{table}[htbp]
\centering
\caption{Bench verification plan: tests, pass marks and the assumption each closes.}
\label{tab:benchplan}
\small
\begin{tabularx}{\textwidth}{lYYl}
\toprule
 & Test & Pass mark & Closes \\
\midrule
B1 & Short circuit on a sacrificial device, hard fault at turn-on, \SI{950}{\volt}, \SI{175}{\celsius}, \SI{+15}{\volt}
   & survives $\geq\SI{0.83}{\micro\second}$; design margin at $\geq\SI{1.66}{\micro\second}$ & A.5 \\
B2 & Parasitic extraction of the power loop, PCB contribution only
   & $\leq\SI{20}{\nano\henry}$ target; C2 lost at \SI{22.6}{\nano\henry} & A.14 \\
B3 & Double pulse on the real board, \SI{950}{\volt}, \SI{60}{\ampere}, hot plate
   & pin voltage $\leq\SI{1055}{\volt}$ (die $\leq\SI{1080}{\volt}$) & A.14 \\
B4 & Switching energies at the datasheet point
   & between the model and the datasheet (\cref{sec:kelvin}) & A.6 \\
B5 & Case-to-plate resistance of the mounted pad, specified torque
   & $\Rthcs \leq \SI{0.18}{\celsius\per\watt}$ design value, \SI{0.50}{\celsius\per\watt} at most (\cref{sec:coolalt}) & A.7 \\
B6 & Locked rotor, \SI{42.4}{\ampere} DC, coolant \SI{65}{\celsius}, case thermocouple on the hottest device
   & $\Tcase \leq \SI{106}{\celsius}$, i.e.\ $\Tj \leq \SI{120}{\celsius}$ at typical $\Rthjc$; model \SI{84.5}{\celsius} on the six-port rail with the aluminium nitride pad (\cref{tab:coolalt}) & C1 \\
B7 & UV inspection of the coating, every board & complete coverage of the HV area & A.1 \\
\bottomrule
\end{tabularx}
\end{table}
